\documentclass{article}
\usepackage[T1]{fontenc}
\usepackage{lmodern}
\usepackage{graphicx}
\usepackage{amsmath,amssymb}
\usepackage[a4paper,margin=2cm]{geometry}
\usepackage[absolute,overlay]{textpos}
\setlength{\TPHorizModule}{1mm}
\setlength{\TPVertModule}{1mm}
\usepackage[indonesian]{babel}
\usepackage{hyperref}
\setlength{\emergencystretch}{2em}
% unit: Halaman 8

\begin{document}

% === Halaman 8 (llm) ===

\begin{itemize}
    \item Pengirim dan penerima pesan memiliki salinan (\textit{copy}) \textit{pad} yang sama.
    \item Satu \textit{pad} hanya digunakan sekali (\textit{one-time}) saja untuk mengenkripsi pesan\\
    $\rightarrow$ itulah mengapa dinamakan \textit{one-time pad}.
    \item Sekali \textit{pad} telah digunakan, ia dihancurkan supaya tidak dipakai kembali untuk mengenkripsi pesan yang lain\\
    $\rightarrow$ menyulitkan kriptanalisis
\end{itemize}

\end{document}
